% Part: turing-machines
% Chapter: machines-computations
% Section: combining-machines

\documentclass[../../../include/open-logic-section]{subfiles}

\begin{document}

\olfileid[pt]{tur}{mac}{cmb}
\olsection{Combinando Máquinas de Turing}

\begin{explain}
Os exemplos de máquinas de Turing que vimos até agora foram bastante
simples por natureza. Mas, de fato, qualquer problema que possa ser resolvido com qualquer
linguagem de programação moderna também pode ser resolvido com máquinas de Turing.
Para construir máquinas de Turing mais complexas, é importante nos convencer
de que podemos combiná-las, de modo a construir máquinas que resolvam
problemas mais complexos quebrando o procedimento em partes mais simples.
Se conseguirmos encontrar uma maneira natural de quebrar um problema complexo em
partes constituintes, podemos atacar o problema em vários estágios,
criando várias máquinas de Turing simples e combinando-as em uma única
máquina capaz de resolver o problema. Esse ponto é especialmente importante
ao atacar o Problema da Parada na próxima seção.

Como combinamos as máquinas de Turing $M = \tuple{Q, \Sigma, q_0, \delta}$
e~$M' = \tuple{Q', \Sigma', q_0', \delta'}$?  Usamos agora a
configuração da fita após $M$~ter parado como a configuração de entrada
de uma execução da máquina~$M'$.  Para obter uma única máquina de Turing
$M \frown M'$ que faz isso, faça o seguinte:
\begin{enumerate}
    \item Renumere (ou rerrotule) todos os estados~$Q'$ de~$M'$ de modo que $M$
    e~$M'$ não tenham estados em comum ($Q \cap Q' = \emptyset$).
    \item Os estados de $M \frown M'$ são $Q \cup Q'$.
    \item O alfabeto da fita é $\Sigma \cup \Sigma'$.
    \item O estado inicial é~$q_0$.
    \item A função de transição é a função $\delta''$ dada por:
    \[\delta''(q,\sigma) =
    \begin{cases}
      \delta(q,\sigma) & \text{se $q \in Q$}\\
      \delta'(q,\sigma) & \text{se $q \in Q'$}\\
      \tuple{q_0', \sigma, \TMstay} & \text{se $q \in Q$ e
      $\delta(q,\sigma)$ é indefinida}
    \end{cases}\]
\end{enumerate}
A máquina resultante usa as instruções de~$M$ quando está em um
estado $q \in Q$, e as instruções de~$M'$ quando está em um estado~$q \in
Q'$. Quando está em um estado $q \in Q$ e examina um símbolo~$\sigma$
para o qual $M$ não tem transição (isto é, $M$ teria parado), ela
entra no estado inicial de~$M'$ (e deixa o conteúdo da fita e a posição
da cabeça como estão).

Note que, a menos que a máquina~$M$ seja disciplinada, não sabemos onde
está a cabeça da fita quando $M$~para, então a configuração de parada de~$M$
não precisa ter a cabeça examinando o quadrado~$1$. Ao combinar máquinas,
é importante ter isso em mente.
\end{explain}

\begin{ex}
\emph{Combinando Máquinas:} Projetaremos uma máquina que, quando iniciada
com uma entrada composta de dois blocos de~$\TMstroke$ de comprimento $n$
e~$m$, para com um único bloco de $2(m+n)$ $\TMstroke$ na
fita. Para construir essa máquina, podemos combinar duas máquinas que
já conhecemos: a máquina de adição e a duplicadora. Começamos
desenhando um diagrama de estado para a máquina de adição.
\[
\begin{tikzpicture}[->,>=stealth',shorten >=1pt,auto,node distance=2.8cm,
                    semithick]
  \tikzstyle{every state}=[fill=none,draw=black,text=black]

  \node[initial,state] (A)              {$q_0$};
  \node[state]         (B) [right of=A] {$q_1$};
  \node[state]         (C) [right of=B] {$q_2$};

  \path (A) edge node {\TMtrans{\TMblank}{\TMstroke}{\TMstay}} (B)
            edge [loop above] node {\TMtrans{\TMstroke}{\TMstroke}{\TMright}} (B)
        (B) edge [loop above] node {\TMtrans{\TMstroke}{\TMstroke}{\TMright}} (B)
            edge node {\TMtrans{\TMblank}{\TMblank}{\TMleft}} (C)
        (C) edge [loop above] node {\TMtrans{\TMstroke}{\TMblank}{\TMstay}} (C);
\end{tikzpicture}
\]
Em vez de parar no estado~$q_2$, queremos continuar a operação a fim
de duplicar a saída. Lembre-se de que a máquina duplicadora apaga o
primeiro traço da entrada e escreve dois traços em uma saída separada.
Vamos adicionar uma instrução para garantir que a cabeça da fita esteja lendo o
primeiro traço da saída da máquina de adição.
\[
\begin{tikzpicture}[->,>=stealth',shorten >=1pt,auto,node distance=2.8cm,
                    semithick]
  \tikzstyle{every state}=[fill=none,draw=black,text=black]

  \node[initial,state] (A)              {$q_0$};
  \node[state]         (B) [right of=A] {$q_1$};
  \node[state]         (C) [right of=B] {$q_2$};
  \node[state]         (D) [below left of=C] {$q_3$};
  \node[state]         (E) [below left of=D] {$q_4$};

  \path (A) edge node {\TMtrans{\TMblank}{\TMstroke}{\TMstay}} (B)
            edge [loop above] node {\TMtrans{\TMstroke}{\TMstroke}{\TMright}} (B)
        (B) edge [loop above] node {\TMtrans{\TMstroke}{\TMstroke}{\TMright}} (B)
            edge node {\TMtrans{\TMblank}{\TMblank}{\TMleft}} (C)
        (C) edge node {\TMtrans{\TMstroke}{\TMblank}{\TMleft}} (D)
        (D) edge [loop left] node {\TMtrans{\TMstroke}{\TMstroke}{\TMleft}} (D)
            edge node[left, xshift=-2mm] {\TMtrans{\TMendtape}{\TMendtape}{\TMright}} (E);
\end{tikzpicture}
\]
Agora é fácil duplicar a entrada---tudo o que temos a fazer é conectar a
máquina duplicadora ao estado~$q_4$. Isso requer renomear os estados da
máquina duplicadora de modo que comecem em~$q_4$ em vez
de~$q_0$---dessa forma não acabamos com dois estados iniciais. O
diagrama final deve ter a aparência da \olref{fig:combined}.
\begin{figure}
\[
\begin{tikzpicture}[->,>=stealth',shorten >=1pt,auto,node distance=2.8cm,
                    semithick]
  \tikzstyle{every state}=[fill=none,draw=black,text=black]
  \node[initial,state] (A)              {$q_0$};
  \node[state]         (B) [right of=A] {$q_1$};
  \node[state]         (C) [right of=B] {$q_2$};
  \node[state]         (D) [below left of=C] {$q_3$};
  \node[state]         (E) [below left of=D] {$q_4$};
  \node[state]         (2) [right of=E] {$q_5$};
  \node[state]         (3) [right of=2] {$q_6$};
  \node[state]         (4) [below of=3] {$q_7$};
  \node[state]         (5) [left of=4]  {$q_8$};
  \node[state]         (6) [left of=5]  {$q_9$};

  \path (A) edge node {\TMtrans{\TMblank}{\TMstroke}{\TMstay}} (B)
            edge [loop above] node {\TMtrans{\TMstroke}{\TMstroke}{\TMright}} (B)
        (B) edge [loop above] node {\TMtrans{\TMstroke}{\TMstroke}{\TMright}} (B)
            edge node {\TMtrans{\TMblank}{\TMblank}{\TMleft}} (C)
        (C) edge node {\TMtrans{\TMstroke}{\TMblank}{\TMleft}} (D)
        (D) edge [loop left] node {\TMtrans{\TMstroke}{\TMstroke}{\TMleft}} (D)
            edge node[left, xshift=-2mm] {\TMtrans{\TMendtape}{\TMendtape}{\TMright}} (E)
    (E) edge node {\TMtrans{\TMstroke}{\TMblank}{\TMright}} (2)
    (2) edge [loop above] node {\TMtrans{\TMstroke}{\TMstroke}{\TMright}} (2)
      edge node {\TMtrans{\TMblank}{\TMblank}{\TMright}} (3)
    (3) edge [loop above] node {\TMtrans{\TMstroke}{\TMstroke}{\TMright}} (3)
        edge node {\TMtrans{\TMblank}{\TMstroke}{\TMright}} (4)
    (4) edge [loop below] node {\TMtrans{\TMblank}{\TMstroke}{\TMleft}} (4)
        edge node {\TMtrans{\TMstroke}{\TMstroke}{\TMleft}} (5)
    (5) edge [loop below]  node {\TMtrans{\TMstroke}{\TMstroke}{\TMleft}} (5)
        edge              node {\TMtrans{\TMblank}{\TMblank}{\TMleft}} (6)
    (6) edge [loop below] node {\TMtrans{\TMstroke}{\TMstroke}{\TMleft}} (6)
        edge              node {\TMtrans{\TMblank}{\TMblank}{\TMright}} (E);
\end{tikzpicture}
\]
\caption{Combinando as máquinas de adição e duplicadora}
\ollabel{fig:combined}
\end{figure}
\end{ex}

\begin{prop}
    Se $M$ e $M'$ são disciplinadas e computam as funções $f\colon
    \Nat^k \to \Nat$ e $f'\colon \Nat \to \Nat$, respectivamente, então
    $M \frown M'$ é disciplinada e computa~$\comp{f}{f'}$.
\end{prop}

\begin{proof}
    Como $M$ é disciplinada, quando para com
    saída~$f(n_1,\dots,n_k) = m$, a cabeça está examinando o quadrado~$1$. Se
    agora entrarmos no estado inicial de~$M'$, a máquina parará com
    saída $f'(m)$, novamente examinando o quadrado~$1$. As demais condições de
    \olref[dis]{defn:disciplined} também são satisfeitas.
\end{proof}

\begin{prob}
    Dê uma máquina de Turing disciplinada que computa $f(x) = x+2$ tomando
    a máquina~$M$ de \cref{tur:mac:dis:prob:disc-succ} e
    construindo $M \frown M$.
\end{prob}
\end{document}
